\chapter{Special Relativity and Tensor Notation}

\section{The invariant interval}

Set \(c=1\).  A spacetime point is
\[
  x^\mu=(x^0,x^1,x^2,x^3)=(t,x,y,z).
\]
The invariant interval is
\[
  s^2=t^2-x^2-y^2-z^2.
\]
Introduce the metric
\[
  \eta_{\mu\nu}=\diag(1,-1,-1,-1).
\]
Then
\[
  s^2=\eta_{\mu\nu}x^\mu x^\nu.
\]
Repeated upper-lower indices are summed:
\[
  a_\mu b^\mu=a_0b^0+a_1b^1+a_2b^2+a_3b^3.
\]
Lowering an index means
\[
  x_\mu=\eta_{\mu\nu}x^\nu=(t,-x,-y,-z).
\]

\section{Lorentz transformations}

A Lorentz transformation is a linear change of coordinates
\[
  x'^\mu=\Lambda^\mu{}_\nu x^\nu
\]
that preserves the interval:
\[
  \eta_{\rho\sigma}x'^\rho x'^\sigma
  =
  \eta_{\mu\nu}x^\mu x^\nu.
\]
Substitute \(x'^\rho=\Lambda^\rho{}_\mu x^\mu\):
\[
  \eta_{\rho\sigma}\Lambda^\rho{}_\mu\Lambda^\sigma{}_\nu x^\mu x^\nu
  =
  \eta_{\mu\nu}x^\mu x^\nu.
\]
Since this holds for all \(x\),
\[
  \eta_{\rho\sigma}\Lambda^\rho{}_\mu\Lambda^\sigma{}_\nu
  =
  \eta_{\mu\nu}.
\]

A boost in the \(x\)-direction is
\[
  t'=\gamma(t-vx),\qquad
  x'=\gamma(x-vt),\qquad
  y'=y,\qquad z'=z,
\]
with
\[
  \gamma=\frac1{\sqrt{1-v^2}}.
\]
Check the interval:
\[
  t'^2-x'^2
  =
  \gamma^2[(t-vx)^2-(x-vt)^2]
  =
  \gamma^2(1-v^2)(t^2-x^2)
  =
  t^2-x^2.
\]

\section{Derivatives and the d'Alembertian}

Define
\[
  \p_\mu=\frac{\p}{\p x^\mu}
  =
  \left(\frac{\p}{\p t},\frac{\p}{\p x},\frac{\p}{\p y},\frac{\p}{\p z}\right).
\]
Raising the index gives
\[
  \p^\mu=\eta^{\mu\nu}\p_\nu
  =
  (\p_t,-\nabla).
\]
The Lorentz-invariant wave operator is
\[
  \Boxop=\p_\mu\p^\mu=\p_t^2-\nabla^2.
\]
The Klein--Gordon equation is
\[
  (\Boxop+m^2)\phi=0.
\]
Its Lagrangian is
\[
  \Lag=\frac12\p_\mu\phi\,\p^\mu\phi-\frac12m^2\phi^2.
\]
Expanding the first term,
\[
  \p_\mu\phi\,\p^\mu\phi=(\p_t\phi)^2-|\nabla\phi|^2.
\]

\section{Four-momentum}

The four-momentum is
\[
  p^\mu=(E,\bm p),\qquad p_\mu=(E,-\bm p).
\]
Its square is
\[
  p_\mu p^\mu=E^2-|\bm p|^2.
\]
For a particle of mass \(m\),
\[
  p^2=m^2.
\]
The plane wave
\[
  e^{-\ii p_\mu x^\mu}=e^{-\ii Et+\ii\bm p\cdot\bm x}
\]
solves the Klein--Gordon equation when
\[
  (-p^2+m^2)e^{-\ii p\cdot x}=0,
\]
so \(p^2=m^2\).

\section{Tensors}

A scalar is unchanged by a Lorentz transformation.  A vector transforms as
\[
  V'^\mu=\Lambda^\mu{}_\nu V^\nu.
\]
A rank-two tensor transforms as
\[
  T'^{\mu\nu}
  =
  \Lambda^\mu{}_\rho\Lambda^\nu{}_\sigma T^{\rho\sigma}.
\]
Contracting one upper with one lower index produces an invariant or a lower-rank tensor.  For example,
\[
  V_\mu W^\mu
\]
is a scalar.

The field strength of electromagnetism will be an antisymmetric tensor \(F_{\mu\nu}=-F_{\nu\mu}\).  The word ``tensor'' matters because it tells us an equation has the same form for all inertial observers.

\section{Natural units and dimensions}

Set \(\hbar=c=1\).  Then time and length have units of inverse energy:
\[
  [x]=[t]=E^{-1},\qquad [\p_\mu]=E.
\]
The action is dimensionless because path integrals contain \(e^{\ii S}\).  Since
\[
  S=\int\dd^4x\,\Lag,
\]
and \([\dd^4x]=E^{-4}\), the Lagrangian density has dimension
\[
  [\Lag]=E^4.
\]
For a scalar field in four spacetime dimensions,
\[
  \Lag\supset \frac12\p_\mu\phi\p^\mu\phi
\]
implies
\[
  E^4=(E[\phi])^2,
  \qquad [\phi]=E.
\]
Dimensional analysis will become a practical tool in effective field theory.

\section*{Exercises}
\addcontentsline{toc}{section}{Exercises}

\begin{exercise}
Show that \(\p_\mu j^\mu=\p_t\rho+\nabla\cdot\bm j\) if \(j^\mu=(\rho,\bm j)\).
\begin{answer}
\(\p_\mu j^\mu=\p_0j^0+\p_ij^i=\p_t\rho+\p_xj_x+\p_yj_y+\p_zj_z\).
\end{answer}
\end{exercise}

\begin{exercise}
Find the mass dimension of a Dirac spinor \(\psi\) in four dimensions from the kinetic term \(\bar\psi\ii\gamma^\mu\p_\mu\psi\).
\begin{answer}
\(E^4=[\bar\psi]E[\psi]\), and \([\bar\psi]=[\psi]\), so \([\psi]=E^{3/2}\).
\end{answer}
\end{exercise}

